\newpage
\chapter*{General Index}
\addcontentsline{toc}{chapter}{General Index}
\label{index:general}
This index lists mainly terms of the \MPI/ specification. The underlined page numbers refer to the definitions or parts of the definition of the terms. Bold face page numbers are used for entries appearing in a section or subsection title.

\include{general-html.tex}

\newpage
\chapter*{Examples Index}
\addcontentsline{toc}{chapter}{Examples Index}

\label{index:example}
This index lists code examples throughout the text.  Some examples are referred to by content; others are listed by the major \MPI/ procedure name that they are demonstrating.  \MPI/ procedure names listed in all capital letter are Fortran examples; \MPI/ procedure names listed in mixed case are C examples.

\include{examples-html.tex}

\newpage
\chapter*{MPI Constant and Predefined Handle Index}
\addcontentsline{toc}{chapter}{MPI Constant and Predefined Handle Index}

\label{index:constant}
This index lists predefined \MPI/ constants and handles, including info keys and values. Underlined page numbers give the location of the primary definition or use of the indexed term.

\include{const-html.tex}

\newpage
\chapter*{MPI Declarations Index}
\addcontentsline{toc}{chapter}{MPI Declarations Index}
\label{index:declaration}
This index refers to declarations needed in C and Fortran, such as address kind integers, handles, etc. The underlined page numbers is the ``main'' reference (sometimes there are more than one when key concepts are discussed in multiple areas). Fortran defined types are shown as \mpicode{TYPE($name$)}.

\include{decls-html.tex}

\newpage
\chapter*{MPI Callback Function Prototype Index}
\addcontentsline{toc}{chapter}{MPI Callback Function Prototype Index}
\label{index:callback}
This index lists the C typedef names for callback routines, such as those used with attribute caching or user-defined reduction operations.  Fortran example prototypes are given near the text of the C name.

\include{typedef-html.tex}

\newpage
\chapter*{MPI Function Index}
\addcontentsline{toc}{chapter}{MPI Function Index}
\label{index:function}
The underlined links refer to the function definitions.

\include{funcdef-html.tex}

